{\small\textbf{Borrowers (Diego, 1:15) [17:40–18:55]}}\par\smallskip
Cheap for banks, but did borrowers pay?
\begin{itemize}\setlength\itemsep{0pt}
\item \textcolor{navy}{\textbf{[def]}} We ran an event study: we compare Dutch lending rates and loan growth with euro-area countries that did not change their buffer, around each of the three policy dates.
\item \textcolor{navy}{\textbf{[def]}} With only one treated country, we judge significance with a placebo test: we pretend each control country raised its buffer and see how large such "fake" effects are. That is the grey band.
\end{itemize}
\par\smallskip
The red line is the Netherlands.
\begin{itemize}\setlength\itemsep{0pt}
\item Its \textbf{lending rates stay inside the band}: no unusual change at any of the three dates.
\item \textbf{Household credit actually grew faster} than in the peer countries.
\end{itemize}
\par\smallskip
One more honest detail. Over 2022 to 2026, Dutch bank credit grew 16.6 percent against 7.9 percent in the euro area. But in the three years before, the Netherlands had lagged, 4.3 against 11.3 percent, so part of that later gap is catching up.
\par\smallskip
Our claim is modest: one treated country and a small control group, so this is "\textbf{no sign of contraction}", not a causal estimate. DNB's own verdict in September 2026 was the same: "no signs of constraints in bank credit supply".
\par\smallskip
\textcolor{nlred}{\textbf{HAND-OVER:}} "So the decision was cheap and credit kept flowing. But a buffer like this reaches beyond Dutch banks and borrowers. Zheng."
